\subsection{\texorpdfstring{测度 \(\mu\) 关于 \(\mu\) -逆紧映射的分解}{测度 mu 关于 mu -逆紧映射的分解}}

设 T 是具有可数基的局部紧空间（换言之，即局部紧的 Polish 空间（GT, IX, §6, No.~1）。我们知道，对 T 上的每个正测度，积分与本质积分这两个概念一致（Ch. V, §1, No.~3, 命题 9 的推论）。另一方面，我们有下列性质：

引理 1. --- 若 Y 是具有可数基的局部紧空间，则空间 \(\mathcal{K}(\mathrm{Y})\) 含有一个可数稠密子集。更确切地说，在 \(\mathcal{K}(\mathrm{Y})\) 中存在一个由 \(\geqslant 0\) 的函数组成的可数子集 S，使得对每个函数 \(f \geqslant 0\) （ \(\mathcal{K}(\mathrm{Y})\) 中的函数），存在一个函数序列 \(f_n \in \mathrm{S}\) （ \(n \geqslant 0\) ），它一致收敛于 \(f\) ，且满足 \(f_n \leqslant f_0\) （对所有 \(n \geqslant 0\) ）。

因为，Y 是相对紧开集的递增序列 \((\mathrm{U}_{n})\) 的并，使得对所有 n 有 \(\overline{U}_{n} \subset U_{n+1}\) （GT, I, §9, No.~9, 命题 15）；空间 \(\mathcal{K}(\mathrm{Y})\) 是 Banach 空间递增序列 \(\mathcal{K}(\mathrm{Y}, \overline{\mathrm{U}}_{n})\) 的并，且已知后者中的每一个都是可分的（GT, Ch. X, §3, No.~3, 定理 1）。令 \(S_{n}^{\prime}\) 为 \(\mathcal{K}(\mathrm{Y}, \overline{\mathrm{U}}_{n})\) 中的可数稠密集， \(S_{n}\) 为诸函数 \(\varphi^{+}\) （ \(\varphi \in S_{n}^{\prime}\) ）所成之集， \(u_{n}\) 为 \(\mathcal{K}(\mathrm{Y}, \overline{\mathrm{U}}_{n+1})\) 中取值于 {[}0,1{]} 且在 \(U_{n}\) 上等于 1 的函数。我们取 S 为诸 \(S_{n}\) 的并以及诸 \(mu_{n}\) （m 与 n 为 \(\geqslant 0\) 的整数）所成之集。对每个函数 \(f \geqslant 0\) （ \(\mathcal{K}(\mathrm{Y})\) 中的函数），f 的支集含于某个 \(U_{n}\) ，故 f 是函数序列 \(f_{p} \in S_{n}\) （ \(p \geqslant 1\) ）的一致极限。这些函数 \(f_{p}\) 被某个正整数 m 一致有界，于是只需取 \(f_{0} = mu_{n}\) 。

引理 2. --- 若 T 是具有可数基的局部紧空间，则由在无穷远点趋于 0 的连续数值函数组成的 Banach 空间 \(\mathcal{K}(\mathrm{Y})\) 是可分的。

这个引理不是别的，正是 GT, X, §3, No.~3 定理 1 的推论。可以注意到，它也可由引理 1 以及如下事实得出： \(\mathcal{H}(\mathrm{Y})\) 上的一致收敛拓扑粗于诸子空间 \(\mathcal{K}(\mathrm{Y},\overline{\mathrm{U}}_n)\) 的拓扑的直接极限拓扑。

引理 3. --- 设 T 与 X 是两个具有可数基的局部紧空间， \(\mu\) 是 T 上的正测度， \(t \mapsto \lambda_{t}\) （ \(t \in T\) ）是 X 上正测度的一个族。若映射 \(t \mapsto \lambda_{t}\) 是标量 \(\mu\) -可积的（对拓扑 \(\sigma(\mathcal{M}(X), \mathcal{K}(X))\) ），则族 \(t \mapsto \lambda_{t}\) 是 \(\mu\) -适宜的（ \(\S1\) ， No.~1, 例）。

因为，把引理 1 应用于 \(\mathcal{K}(\mathrm{X})\) ，便表明映射 \(t \mapsto \lambda_t\) 是淡 \(\mu\) -可测的（§1, No.~5, 命题 13）。

定理 1. --- 设 T 与 B 是两个具有可数基的局部紧空间， \(\mu\) 是 T 上的正测度， \(p\) 是 T 到 B 中的 \(\mu\) -逆紧映射（Ch. V, §6, No.~1, 定义 1）， \(\nu = p(\mu)\) 是 \(\mu\) 在 \(p\) 下的像。则存在 T 上正测度的一个 \(\nu\) -适宜族（§1, No.~1, 例） \(b \mapsto \lambda_b\) （ \(b \in B\) ），它具有下列性质：

\begin{enumerate}
\def\labelenumi{\alph{enumi})}
\item
  \(\| \lambda_b\| = 1\) （对一切 \(b\in p(\mathrm{T})\) ）
\item
  \(\lambda_{b}\) 集中于集合 \(\overline{p}^{-1}(b)\) 上（Ch. V, §5, No.~7, 定义 4）（对一切 \(b\in \mathrm{B}\) ）；特别地， \(\lambda_{b} = 0\) （当 \(b\notin p(\mathrm{T})\) 时）；
\item
  \(\mu = \int \lambda_b d\nu(b)\)。
\end{enumerate}

此外，若 \(b \mapsto \lambda_b'\) （ \(b \in \mathbf{B}\) ）是另一个 \(\nu\) -适宜族（ \(\mathbf{T}\) 上的正测度族），具有性质 \(\mathbf{b}\) ）与 \(\mathbf{c}\) ），则 \(\lambda_b' = \lambda_b\) 在 \(\mathbf{B}\) 中关于测度 \(\nu\) 几乎处处成立。

\begin{enumerate}
\def\labelenumi{\arabic{enumi})}
\tightlist
\item
  唯一性。对每个函数 \(f \in \mathcal{K}(\mathrm{B})\) ， \(f \circ p\) 是 \(\mu\) -可积的，因为 \(p\) 是 \(\mu\) -逆紧的（Ch. V, §6, No.~2, 定理 1）；于是对每个函数 \(g \in \mathcal{K}(\mathrm{T})\) ，函数 \(t \mapsto g(t)f(p(t))\) 是 \(\mu\) -可积的。由此得出（Ch. V, §3, No.~3, 定理 1）：对几乎每个 \(b \in \mathrm{B}\) ，函数 \(t \mapsto g(t)f(p(t))\) 是 \(\lambda_b\) -可积的，且
\end{enumerate}

\[
\int g (t) f \bigl (p (t) \bigr) d \mu (t) = \int d \nu (b) \int g (t) f \bigl (p (t) \bigr) d \lambda_ {b} (t).\tag{1}
\]

但由于 \(\lambda_{b}\) 集中于 \(\overline{p}^{-1}(b)\) 上，故对每个 \(b \in \mathbf{B}\) ， \(f(p(t)) = f(b)\) 关于 \(\lambda_{b}\) 几乎处处成立，因此 (1) 的第二端等于 \(\int f(b)\langle g, \lambda_{b} \rangle d\nu(b)\) 。关于 \(\lambda_{b}'\) 的类似公式也成立；于是 \(\int f(b)\langle g, \lambda_{b} \rangle d\nu(b) = \int f(b)\langle g, \lambda_{b}' \rangle d\nu(b)\) （对一切 \(f \in \mathcal{K}(\mathbf{B})\) 与 \(g \in \mathcal{K}(\mathbf{T})\) ）。换言之，两个映射 \(b \mapsto \lambda_{b}\) 与 \(b \mapsto \lambda_{b}'\) （由 B 到 \(\mathcal{M}(T)\) 中的映射）关于 \(\nu\) 标量意义局部几乎处处相等，从而关于 \(\nu\) 几乎处处相等（引理 1 与 §1, No.~1, 注记 2）。

\begin{enumerate}
\def\labelenumi{\arabic{enumi})}
\setcounter{enumi}{1}
\tightlist
\item
  族 \(b \mapsto \lambda_b\) 的临时定义。对每个函数 \(f \in \mathcal{L}^1(\nu)\) ， \(f \circ p\) 是 \(\mu\) -可积的（Ch. V, §6, No.~2, 定理 1），因此 \((f \circ p) \cdot \mu\) 是 \(T\) 上的有界测度，且
\end{enumerate}

\[
\left\| (f \circ p) \cdot \mu \right\| = \int | f \circ p | d \mu = \int | f | d \nu = \mathrm{N} _ {1} (f)
\]

（Ch. IV, §4, No.~7, 命题 12；Ch. V, §5, No.~3, 定理 1 与 §6, No.~2, 定理 1）。由此可知， \((f \circ p) \cdot \mu\) 只依赖于类 \(\widetilde{f}\) （ \(f\) 在 \(\mathrm{L}^1(\nu)\) 中的类），且 \(\widetilde{f} \mapsto (f \circ p) \cdot \mu\) 是 \(\mathrm{L}^1(\nu)\) 到 Banach 空间 \(\mathcal{M}^1(\mathrm{T})\) （ \(\mathrm{T}\) 上有界测度所成的 Banach 空间，即 Banach 空间 \(\overline{\mathcal{K}(\mathrm{T})}\) 的强对偶，后者可分（引理 2））中的等距线性映射。根据 Dunford-Pettis 定理（§2, No.~5, 定理 1 的推论 2），存在一个映射 \(b \mapsto \lambda_b\) ，由 \(\mathrm{B}\) 取值于 \(\mathcal{M}^1(\mathrm{T})\) 的单位球中，标量 \(\nu\) -可测（对拓扑 \(\sigma(\mathcal{M}^1(\mathrm{T}), \overline{\mathcal{K}(\mathrm{T})})\) ），且使得对每个函数 \(f \in \mathcal{L}^1(\nu)\) ，

\[
(f \circ p) \cdot \mu = \int f (b) \lambda_ {b} d \nu (b),\tag{2}
\]

对每个函数 \(g \in \overline{\mathcal{K}(\mathrm{T})}\) ，它也可写成

\[
\int g (t) f (p (t)) d \mu (t) = \int f (b) d \nu (b) \int g (t) d \lambda_ {b} (t).\tag{3}
\]

若 \(f \geqslant 0\) 且 \(g \geqslant 0\) ，则 (3) 的第一端 \(\geqslant 0\) ，这就证明了对每个函数 \(g \geqslant 0\) （ \(\mathcal{K}(\mathrm{T})\) 中的函数），测度 \(\left( \int g(t) d\lambda_b(t) \right) \cdot \nu\) \(\geqslant 0\) ，从而 \(\int g(t) d\lambda_b(t) \geqslant 0\) （ \(b\) 属于一个 \(\nu\) -可忽略集 \(\mathrm{N}(g)\) 者除外）（Ch. V, §5, No.~3, 命题 3 的推论 3）。现在，存在一个稠密序列 \((g_n)\) ，位于空间 \(\mathcal{K}_+(\mathrm{T})\) （由 \(\geqslant 0\) 且属于 \(\mathcal{K}(\mathrm{T})\) 的函数所成的空间）中（引理 1）。并 \(\mathrm{N}\) （即诸 \(\mathrm{N}(g_n)\) 的并）是 \(\nu\) -可忽略的，且对 \(b \notin \mathrm{N}\) ，有 \(\int g_n(t) d\lambda_b(t) \geqslant 0\) （对一切 \(n\) ），因此 \(\int g(t) d\lambda_b(t) \geqslant 0\) （对每个函数 \(g \in \mathcal{K}_+(\mathrm{T})\) ），换言之 \(\lambda_b \geqslant 0\) 。

这样一来，我们可以把 \(\lambda_{b}\) （对每个 \(b \in N\) ）换成 0 而不改变 (3) 的有效性；因此可以假定这一修改已经完成，从而 \(\lambda_{b} \geqslant 0\) （对每个 \(b \in B\) ）。

\begin{enumerate}
\def\labelenumi{\arabic{enumi})}
\setcounter{enumi}{2}
\tightlist
\item
  公式 (3) 的推广。
\end{enumerate}

\(\alpha)\) 对每个函数 \(f \in \mathcal{L}^1(\nu)\) ，由 (3) 推出，映射 \(b \mapsto \lambda_b\) （ B 到 \(\mathcal{M}(T)\) 中的映射）关于测度 \(|f \cdot \nu|\) 与拓扑 \(\sigma(\mathcal{M}(T), \mathcal{K}(T))\) 是标量可积的，因此（引理 3）族 \(b \mapsto \lambda_b\) 是 \(|f \cdot \nu|\) -适宜的。现在设 \(g\) 是一个数值函数（定义在 \(T\) 上），关于测度 \(|(f \circ p) \cdot \mu|\) 可积，也就是说（Ch. V, §5, No.~3, 定理 1）， \(t \mapsto g(t)f((p(t))\) 是 \(\mu\) -可积的；于是由 (2)、Ch. V, §3, No.~3 的定理 1 与 Ch. V, §5, No.~3 的定理 1 推出：对几乎每个 \(b \in \mathbf{B}\) ， \(g\) 关于 \(\lambda_b\) 可积，（几乎处处有定义的）函数 \(b \mapsto \int g(t)d\lambda_b(t)\) 关于 \(|f \cdot \nu|\) 可积，且公式 (3) 仍然成立。

\(\beta)\) 对每个函数 \(g \in \overline{\mathcal{K}(\mathrm{T})}\) ，把 (3) 应用于 \(f \in \mathcal{K}(\mathrm{B})\) ，便推出映射 \(p\) 关于测度 \(|g \cdot \mu|\) 是逆紧的（Ch. V, §6, No.~1, 定义 1），且在 \(p\) 下测度 \(g \cdot \mu\) 的像是以 \(b \mapsto \int g(t) d\lambda_b(t)\) 为密度（关于 \(\nu\) ）的测度。若随后取 \(f\) 为使 \(f \circ p\) 关于测度 \(|g \cdot \mu|\) 可积的函数，也就是使 \(t \mapsto g(t)f(p(t))\) 为 \(\mu\) -可积的函数（Ch. V, §5, No.~3, 定理 1），则公式 (3) 仍然成立（Ch. V, §6, No.~2, 定理 1）。

\begin{enumerate}
\def\labelenumi{\arabic{enumi})}
\setcounter{enumi}{3}
\tightlist
\item
  族 \(b \mapsto \lambda_b\) 的性质。由性质 \(\beta\) ），我们可以取 \(f = 1\) 、 \(g \in \mathcal{K}(\mathrm{T})\) 来应用公式 (3)；这就证明 \(b \mapsto \lambda_b\) 是标量 \(\nu\) -可积的（对拓扑 \(\sigma(\mathcal{M}(\mathrm{T}), \mathcal{K}(\mathrm{T}))\) ），从而是 \(\nu\) -适宜的（引理 3），且 \(\mu = \int \lambda_b d\nu(b)\) 。
\end{enumerate}

现在设 \(\psi\) 是 \(\mathcal{K}(\mathrm{B})\) 中任意一个函数；性质 \(\alpha\) 的条件可通过取 \(f \in \mathcal{K}(\mathrm{B})\) 与 \(g = \psi \circ p\) 而得到满足，因为函数 \(\psi(p(t))f(p(t))\) 是 \(\mu\) -可积的——由于 \(f\psi\) 属于 \(\mathcal{K}(\mathrm{B})\) 且 \(p\) 是 \(\mu\) -逆紧的。于是 \(\psi \circ p\) 是 \(\lambda_b\) -可积的（对几乎每个 \(b \in \mathrm{B}\) ），且

\[
\int f (p (t)) \psi (p (t)) d \mu (t) = \int f (b) d \nu (b) \int \psi (p (t)) d \lambda_ {b} (t);
\]

但按定义第一端就是 \(\int f(b)\psi (b)d\nu (b)\) 。于是我们看到，对每个函数 \(\psi \in \mathcal{K}(\mathrm{B})\) ，测度 \(\psi \cdot \nu\) 与以 \(b\mapsto \int \psi (p(t))d\lambda_b(t)\) 为密度的测度相同。因此（Ch. V, §5, No.~3, 命题 3 的推论 2）存在一个 \(\nu\) -可忽略集 \(\mathrm{N}'(\psi)\) ，使得对每个 \(b\notin \mathrm{N}'(\psi)\) ，函数 \(\psi \circ p\) 是 \(\lambda_{b}\) -可积的，且 \(\psi (b) = \int \psi (p(t))d\lambda_b(t)\) 。

设 S 是 \(\mathcal{K}(\mathrm{B})\) 的一个可数子集，具有引理 1 中所述的性质（取 Y = B），又设 N' 是 \(\nu\) -可忽略集，即诸 N'(ψ) （ \(\psi \in \mathrm{S}\) ）的并。每个函数 \(\psi \geqslant 0\) （ \(\mathcal{K}(\mathrm{B})\) 中的函数）都是 S 中元素组成的一个序列 \((\psi_n)\) （ \(\psi_n \leqslant \psi_0\) ）的一致极限。于是对 \(b \notin \mathrm{N}'\) ，Lebesgue 定理一方面表明 \(\psi \circ p\) 是 \(\lambda_b\) -可积的，换言之 \(p\) 是 \(\lambda_b\) -逆紧的，另一方面表明 \(\psi(b) = \int \psi(p(t)) d\lambda_b(t)\) 。换言之，映射 \(b \mapsto \varepsilon_b\) 与 \(b \mapsto p(\lambda_b)\) （ B 到 \(\mathcal{M}(\mathrm{B})\) 中的映射，后者几乎处处有定义）关于 \(\nu\) （且对拓扑 \(\sigma(\mathcal{M}(\mathrm{B}), \mathcal{K}(\mathrm{B}))\) ）标量意义几乎处处相等；由此推出这些映射关于 \(\nu\) 几乎处处相等（引理 1 与 §1, No.~1, 注记 2）。最后，若 \(p(\lambda_b) = \varepsilon_b\) ，则集合 B - \{b\} 是 \(\varepsilon_b\) -可忽略的，因此集合 T - \(\overline{p}^{-1}(b)\) 是 \(\lambda_b\) -可忽略的（Ch. V, §6, No.~2, 命题 2 的推论 2），换言之 \(\lambda_b\) 集中于 \(\overline{p}^{-1} (b)\) 上；而另一方面， \(\| \lambda_b\| = \int d\lambda_b = \int d((p(\lambda_b)) = \| \varepsilon_b\| = 1\) （Ch. V, §6, No.~2, 定理 1）。

\begin{enumerate}
\def\labelenumi{\arabic{enumi})}
\setcounter{enumi}{4}
\tightlist
\item
  族 \(b \mapsto \lambda_{b}\) 的修改。这样我们就定义了 T 上测度的一个 \(\nu\) -适宜族 \(b \mapsto \lambda_{b}\) （ \(\geqslant 0\) 的测度），它满足叙述中的条件 c)，且使得对几乎每个 \(b \in B\) ， p 是 \(\lambda_{b}\) -逆紧的， \(\lambda_{b}\) 集中于 \(\overline{p}^{-1}(b)\) 上且范数为 1。设 \(N^{\prime\prime}\) 是 \(\nu\) -可忽略集，由使上述最后三个性质之一不成立的点 \(b \in B\) 组成；于是我们可以按下述方式修改 \(\lambda_{b}\) 。若 \(b \in B - p(T)\) ，取 \(\lambda_{b} = 0\) ；若 \(b \in p(T) \cap N^{\prime\prime}\) ，取 \(\lambda_{b} = \varepsilon_{\xi(b)}\) ，其中 \(\xi(b)\) 是 \(\overline{p}^{-1}(b)\) 中任意一点。由于 \(B - p(T)\) 是 \(\nu\) -可忽略的（Ch. V, §6, No.~2, 命题 2 的推论 3），我们只是在一个可忽略集的点上修改了 \(\lambda_{b}\) ，因而族 \(b \mapsto \lambda_{b}\) 仍是 \(\nu\) -适宜的且具有性质 c)；此外它现在还满足 a) 与 b)，证明至此完成。
\end{enumerate}

T 上正测度的每一个 \(\nu\) -适宜族 \(b \mapsto \lambda_{b}\) ，若具有定理 1 的性质 b) 与 c)，就称为测度 \(\mu\) 关于 \(\mu\) -逆紧映射 p 的一个分解。

